%% Este archivo 'anafa-template.tex' es una plantilla para producir documentos camera-ready para la revista Anales de la Asociación Física Argentina    
%% https://anales.fisica.org.ar/
%% Versión actualizada el 13/02/2020 por Mónica B. Agüero

\documentclass{analesafa}
% Incluir aquí los paquetes adicionales que el autor desee usar:

%===============================================================
% TÍTULO DEL TRABAJO (edición: AUTORES)
%1er campo en castellano, 2do campo en inglés
\titulo{PREPARACIÓN DE MANUSCRITOS “CAMERA-READY” PARA LOS ANALES AFA}
{PREPARATION OF CAMERA-READY MANUSCRIPTS FOR ANALES AFA}
%--------------------------------------------------------------
% AUTORES (edición: AUTORES)
\autor{1}{X.X. Alfa}
\autor{2}{Y.Y. Beta}
\corresponding{1,2}{Z.Z. Gamma} % autor de contacto
%--------------------------------------------------------------
% FILIACIONES DE LOS AUTORES (edición: AUTORES)
\filiacion{1}{Instituto de Matemática Aplicada San Luis (IMASL) – Universidad Nacional De San Luis – CONICET \par Ejército de los Andes 950 – (5700) San Luis – Argentina}
\filiacion{2}{Centro de Investigaciones en Láseres y Aplicaciones (CEILAP), UNIDEF (CITEDEF - CONICET), \par Juan Bautista de La Salle 4397 (B1063ALO), Villa Martelli, Buenos Aires, Argentina.}
%-------------------------------------------------------------
% FECHA (edición: ANALES AFA)
\fecha{Recibido: xx/xx/xxxx; Aceptado: xx/xx/xxxx} 
\setcounter{Maxaffil}{0}
%---------------------------------------------------------------
% Inicio del documento
\begin{document}
%---------------------------------------------------------------
% RESUMEN (edición: AUTORES. Ir a abstract.tex)
\input{abstract}
%---------------------------------------------------------------
% CORREO ELECTRÓNICO DEL AUTOR DE CONTACTO (edición: AUTORES)
\email{* zgamma@citedef.gob.ar}  
%---------------------------------------------------------------
% FORMATO 1er PÁGINA (edición: ANALES AFA)
\thispagestyle{empty}     
%\thispagestyle{firstpage}  

%===============================================================
% CUERPO PRINCIPAL DEL TRABAJO (edición: AUTORES)
\section{INTRODUCCIÓN}

Estas instrucciones tienen como objetivo ayudar a los autores que deseen publicar en Anales AFA a preparar sus manuscritos. Una vez que la versión final del trabajo haya sido aceptada, el autor recibirá la \textit{Prueba de Galera} para su revisión y aceptación final.

Este documento es en sí mismo un ejemplo de cómo quedarán los trabajos en su forma final para publicación. Los artículos pueden ser escritos en ESPAÑOL, PORTUGUÉS o INGLÉS. El título, resumen y palabras clave deberán estar escritos tanto en el idioma elegido como en un segundo idioma. Esta pauta tiene el propósito de cumplir con los requisitos de indexación de la revista en distintos catálogos.
Si elije escribir el artículo en español o portugués, deberá incluir el título, resumen y palabras clave tanto en el idioma elegido y como en INGLÉS; esto último es a los fines de su inclusión en los sistemas internacionales de citaciones. 

Se solicita que todos y cada uno de los autores proporcionen su número de identificador normalizado ORCID; para mayor información, ingresar a \urlstyle{same}\url{https://orcid.org/}.


\section{GUÍA PARA LOS AUTORES}
\subsection{Instrucciones para la escritura}

Los manuscritos deben inscribirse en un rectángulo de $17 \times \unit[25,2]{cm^{2}}$, como se muestra en la Fig. \ref{fig01}, en papel formato DIN A4 ($21 \times \unit[29,7]{cm^{2}}$) y respetando los márgenes que se indican. Las ecuaciones, tablas, figuras y referencias deben seguir un esquema numérico secuencial.
% EJEMPLO: FIGURA EN 1 COLUMNA
 \begin{figure}[ht]
    \centering
    \includegraphics[width=0.44\textwidth]{fig1.jpg}
  \caption{Formato de página recomendado}
  \label{fig01}
  \end{figure}
Los manuscritos deben enviarse en formato PDF para la revisión por pares. Nuestro sistema de envío de manuscritos (visitar  \urlstyle{same}\url{https://anales.fisica.org.ar/}) le guiará paso a paso en el proceso para ingresar los datos de su artículo y cargar los archivos. El autor de correspondencia recibirá información sobre el proceso de revisión de su manuscrito a través del email.  
 
 % EJEMPLO: TABLA QUE ABARCA 2 COLUMNAS
\begin{table*}[ht]
\caption{Estilos recomendados.}
\label{tab_01}
    \centering
    \small
    \begin{tabular}{|m{2cm}|m{1.5cm}|m{7.8cm}|c|}
    \hline
    \centering
    \textbf{Estilo} & \centering\textbf{Tamaño de la fuente} & \centering\textbf{Descripción} & \textbf{Ejemplo} \\ 
    \thickhline
    \centering Título General & \centering 14 pt. & Fuente: negrita; sin sangría; párrafo centrado; interlineado sencillo; espaciado posterior: 12 pt.  &  \large\textbf{TÍTULO} \\ 
    \hline
    \centering Autor & \centering 11 pt. &  Fuente: negrita; sin sangría; párrafo centrado; interlineado sencillo; espaciado posterior: 12 pt. & \large\textbf{{Autor}} \\ 
    \hline
    \centering Dirección & \centering 9 pt. & Fuente: itálica; sin sangría; párrafo centrado; interlineado sencillo; espaciado anterior y posterior: 0 pt.  &  \small\textit{{Dirección}}\\
    \hline
    \centering Revisión y aceptación & \centering 9 pt.  & Fuente: normal; párrafo centrado; interlineado sencillo; espacio entre caracteres expandido a 1 pt.; espaciado anterior: 12 pt; espaciado posterior: 6 pt.  & \small{Aceptado}\\
    \hline
    \centering Resumen & \centering 9 pt. & Fuente: normal; indentado a derecha e izquierda: 1 cm; párrafo centrado; interlineado sencillo; espaciado anterior: 0 pt.; espaciado posterior: 6 pt.  & \small{Resumen}  \\
    \hline
    \centering Palabras clave  &  \centering 9 pt.  & Fuente: itálica; sin sangría; justificado a izquierda; interlineado sencillo; espaciado anterior y posterior: 6 pt.  & \small{\textit{Palabras clave}}  \\
    \hline
    \centering Texto  & \centering 10 pt.  & Fuente: normal; sangría de primera línea: 0,5 cm; justificado a derecha e izquierda; interlineado sencillo; espaciado anterior y posterior: 0 pt.  & Este es un ejemplo ...  \\
    \hline
    \centering Título 1  & \centering 10 pt. &  Fuente: negrita; sin sangría; justificado a izquierda; interlineado sencillo; espaciado anterior: 6 pt.; espaciado posterior: 3 pt.   &  \textbf{I. SECCIÓN} \\
    \hline
    \centering Título 2 & \centering 10 pt.  & Fuente: negrita; sin sangría; justificado a izquierda; interlineado sencillo; espaciado anterior: 6 pt.; espaciado posterior: 3 pt.  & \textbf{Subsección}  \\
    \hline
    \centering Ecuaciones & \centering 10 pt.  &  Centrado en la columna; puede incluir número de ecuación justificado a la derecha; espaciado anterior: 0 pt.; espaciado posterior: 3 pt.  &  $p= 2\pi d$ \hspace{0.5cm}  (1)  \\
    \hline
    \centering Título de tablas  & \centering 10 pt.  & Fuente: itálica; sin sangría; párrafo centrado; interlineado sencillo; espaciado anterior: 6 pt.; espaciado posterior: 9 pt.  &  \textit{TABLA 1}  \\
    \hline
    \centering Descripción de figuras  & \centering 9 pt.  &  Fuente: itálica; sin sangría; justificado a izquierda; interlineado sencillo; espaciado anterior: 6 pt.; espaciado posterior: 9 pt.  &  \textit{Figura 1:}  \\
    \hline
    \centering Agradecimientos & \centering 10 pt. &  Fuente: negrita; sin sangría; justificado a izquierda; interlineado sencillo; espaciado anterior: 6 pt.; espaciado posterior: 3 pt. No numerar  esta sección y situarla al final del manuscrito (antes de las Referencias). &  \textbf{AGRADECIMIENTOS}\\
    \hline
    \centering Referencias & \centering 9 pt. & Sangría francesa: 0,5 cm; justificado a derecha e izquierda; interlineado sencillo; espaciado anterior: 0 pt.; espaciado posterior: 3 pt. No numerar esta sección.   &    \\
    \hline
    \centering Nota al pie  &  \centering 9 pt.  & Fuente: itálica; sin sangría; justificado a izquierda; interlineado sencillo; espaciado anterior y posterior: 0 pt.  &  \small{\textit{* email}}   \\
    \hline
    \end{tabular}
\end{table*}

\normalsize

La versión final, ya corregida y aprobada, se entregará tanto en PDF como en una versión editable. Para este último se aceptarán archivos en Word (extensión .doc o .docx) o \LaTeX{ }  (archivos del manuscrito .tex y, en caso que corresponda, archivo .bib). Además, se solicita a los autores adjuntar las correspondientes figuras en formato JPG ó PNG en ambos casos. Los archivos deberán ser enviados por correo electrónico a la dirección \textcolor{blue}{analesafaba@gmail.com}  y con copia a la dirección \textcolor{blue}{analesafatan@gmail.com}.

El presente documento, así como las instrucciones en general y sobre cómo publicar en Anales AFA se encuentra disponible en la página de la revista. Este documento incluye las especificaciones de todos los estilos necesarios; puede por tanto ser utilizado como plantilla sobreescribiendo en el mismo. Por favor, NO INCLUYA numeración de páginas.

\subsection{Tipos de letra y espaciado de líneas}

Todo manuscrito debe tener espaciado simple, incrementándolo automáticamente para subíndices y supraíndices. Debe usarse fuente Times New Roman y generarse en dos columnas. Se ahorran todos estos ajustes al usar esta plantilla en \LaTeX.
En la Tabla \ref{tab_01} se muestran los tamaños de fuentes, formatos de párrafo y separación entre secciones. Por ejemplo, el párrafo normal requiere una sangría de \unit[1]{cm} y una separación de 6 pts. entre párrafos. Recuerde que 1 \emph{point} = 1 pt. $\cong$ \unit[0.25]{mm}. La lista de Estilos de la Tabla \ref{tab_01} también está incluida en el archivo asociado al presente documento.

\subsection{Tablas}

Las tablas deben colocarse lo más cerca posible del punto donde se hace referencia por primera vez y deben citarse consecutivamente en orden numérico usando números arábigos (1, 2, 3, etc.). 

Los títulos de las tablas deben aparecer directamente arriba de las mismas.  Para generar una línea horizontal gruesa en la tabla usar la opción \color{Blue}\textsf{$\backslash$thickhline}\color{Black}. 

%EJEMPLO: TABLA QUE ABARCA 1 COLUMNA
\begin{table}[ht]
\caption{ Opciones del paquete \textsf{array} (entorno \textsf{tabular}).}
\label{tab_02}
    \centering
    \small           % modificación del tamaño de la fuente
    \begin{tabular}{|m{5.5cm}|c|}
    \hline
    \centering
    \textbf{Descripción} & \textbf{Símbolo}  \\ 
    \thickhline
    Alineación a la izquierda & l  \\ 
    \hline
    Alineación centrada & c  \\ 
    \hline
    Alineación a la derecha & r \\
    \hline
    Genera una caja del ancho indicado & p\{width\} \\
    \hline
    Inserta una línea vertical entre columnas & $\mid$ \\
    \hline
    Genera una caja del ancho indicado. La entrada se centrará verticalmente.  &  m\{width\}  \\
    \hline
    \end{tabular}
\end{table}

\normalsize     % fuente en tamaño normal

El ancho de las tablas puede extenderse a las dos columnas. Para ello se recomienda usar la orden \color{Blue} \textsf{$\backslash$begin\{table*\} ... $\backslash$end\{table*\}} \color{Black}. En esta plantilla se presentan ejemplos de cada caso (ver Tablas \ref{tab_01} y \ref{tab_02}).


\subsection{Figuras}

Las figuras se incluyen en el documento en el sitio en que el autor lo considere más conveniente. Numere cada figura consecutivamente con un número arábigo de acuerdo con el orden en que se discute en el documento. Las leyendas de las figuras deben anotarse debajo de ellas (ver Fig. \ref{fig04}). 

\begin{figure}[ht]
\centering
   \includegraphics[width=0.48\textwidth]{fig2.jpg}
\caption{Modelo de figura en una columna.}
\label{fig04}
\end{figure} 

El ancho de las figuras también puede extenderse a las dos columnas (ver Fig. \ref{fig05}).

Trate los gráficos de modo que sean completamente autoexplicativos. Rotule cada eje (horizontal y vertical) incluyendo las correspondientes unidades, que deben estar separadas y encerradas entre corchetes (ejemplo: \textit{I (mA)}). Si es posible no use potencias de 10; en su lugar, use los prefijos apropiados del Sistema Internacional (ver Tabla \ref{tab_03}).
% EJEMPLO: FIGURA EN DOS COLUMNAS
\begin{figure*}[ht]
\centering
   \includegraphics[width=0.9\textwidth]{fig3.jpg}
\caption{Modelo de figura ocupando dos columnas.}
\label{fig05}
\end{figure*} 


\subsection{Formato de las ecuaciones}

Las ecuaciones deben centrarse en la columna y con su número de orden entre paréntesis, alineado por derecha. La numeración debe ser consecutiva usando números arábigos. Se recomienda utilizar los entornos de ecuaciones disponibles en el paquete \textcolor{Blue}{\textsf{amsmath}},  recordando que
la orden \textcolor{Blue}{\textsf{$\backslash$begin$\{$equation$\}$}} numera automáticamente las ecuaciones. Para escribir ecuaciones sin numeración utilice  \textcolor{Blue}{\textsf{$\backslash$begin$\{$equation*$\}$}}.
Trate el texto y las fórmulas matemáticas como una entidad. Puntúe las expresiones matemáticas como oraciones o partes de oraciones. 
% ECUACIONES EN UNA LÍNEA 
Ejemplo: La solución es de la forma
\begin{equation}
r_{s} = A e^{\alpha^{+}s} + A e^{\alpha^{-}s}
\label{eq01}
\end{equation}
donde 
\begin{equation*}
e^{\alpha^{\pm}} = b \pm \sqrt{b^{2}+1}.
\label{eq01a}
\end{equation*}


Las expresiones matemáticas a menudo deben mostrarse en dos o más líneas debido a limitaciones en el ancho de la columna definido por la revista. Las ecuaciones (\ref{eq02}) y (\ref{eq03}) son ejemplos de este caso. Para estos casos se recomienda utilizar los entornos  \\
\color{Blue}
\textsf{$\backslash$begin$\{$eqnarray$\}$} \\ \textsf{$\backslash$end$\{$eqnarray$\}$} \\
\color{black} 

% ECUACIÓN EN DOS LÍNEAS 
\begin{eqnarray}
\beta(\mathbf{k}) & \approx & \beta^{(0)}\,\delta(\mathbf{k}-\mathbf{k}_i) + \beta^{(1)}\,Z(\mathbf{k}) \nonumber\\
 & + & \int d^2k'\,\beta^{(2)}(\mathbf{k},\mathbf{k}')\,Z(\mathbf{k}-\mathbf{k}'),
\label{eq02}
\end{eqnarray}
para generar una ecuación en dos renglones ó  bien \\
\color{Blue}
\textsf{$\backslash$begin$\{$equation$\}$} \\ 
\textsf{$\backslash$begin$\{$split$\}$}  \\
\textsf{$\backslash$end$\{$split$\}$}  \\
\textsf{$\backslash$end$\{$equation$\}$}
\color{black}


% ECUACIÓN EN DOS LÍNEAS 
\begin{equation}
\begin{split}
& E_{x}=-\frac{1}{c}\frac{\partial A_{x}}{\partial t}-\frac{\partial \Phi}{\partial x} = - (\partial^{0}A^{1}-\partial^{1}A^{0}) \\
& B_{x}=\frac{\partial A_{z}}{\partial y}-\frac{\partial A_{y}}{\partial z}=-(\partial^{2}A^{3}-\partial^{3}A^{2}),
\end{split}
\label{eq03}
\end{equation}
para producir dos (o más) ecuaciones etiquetadas con un sólo número.


Si la ecuación es muy larga, como en la ecuación (\ref{eq04}), y no es posible o apropiado mostrarla en dos o más líneas, puede optar por presentarla en dos columnas. En este caso puede emplear un flotante extendido con el comando \textcolor{Blue}{\textsf{$\backslash$begin$\{$figure*$\}$}}
para posicionar la ecuación en el ancho total de la página.

% ECUACIÓN EN DOS COLUMNAS
\begin{figure*}
\begin{eqnarray}
[\mathbf{E}(\mathbf{\nabla . E}) - \mathbf{E}\times (\mathbf{\nabla \times E})]_{1}  & = & 
E_{1}\left( \frac{\partial E_{1}}{\partial x_{1}} +  \frac{\partial E_{2}}{\partial x_{2}} + \frac{\partial E_{3}}{\partial x_{3}} \right) 
- E_{2}\left( \frac{\partial E_{2}}{\partial x_{1}} -  \frac{\partial E_{1}}{\partial x_{2}} \right)
+ E_{3}\left( \frac{\partial E_{1}}{\partial x_{3}} -  \frac{\partial E_{3}}{\partial x_{1}} \right)
\nonumber\\
 & = & \frac{\partial}{\partial x_{1}}(E_{1}^{2}) + \frac{\partial}{\partial x_{2}}(E_{1}E_{2}) + \frac{\partial}{\partial x_{3}}(E_{1}E_{3}) - \frac{1}{2}\frac{\partial}{\partial x_{1}}(E_{1}^{2}+E_{2}^{2}+E_{3}^{2})
\label{eq04}
\end{eqnarray}
\end{figure*}

\subsection{Agrupación de símbolos: jerarquía}

Con el propósito de agrupar, la secuencia de corchetes preferida para los artículos publicados en Anales AFA es \large{\{[()]\}}. \normalsize Ejemplo:
\begin{equation}
D=3\left[L^{2}+\frac{R}{(L+1.5)^{2}}\right]+(R-3L)R. 
\label{eq05}
\end{equation}

\subsection{Unidades}

Utilice las reglas y convenciones aceptadas internacionalmente, como el Sistema Internacional de unidades (SI). Para acceder a una lista de símbolos estándar para unidades de medida  puede  consultar en \urlstyle{same}\url{https://doi.org/10.6028/NIST.SP.330-2019}.

Los múltiplos decimales o submúltiplos de unidades se indican mediante el uso de prefijos como se muestra en la Tabla \ref{tab_03}).

\begin{table}[ht]
\caption{Prefijos del Sistema Internacional. }
\label{tab_03}
    \centering
    \small               % modificación del tamaño de la fuente
    \begin{tabular}{|c|c|c|c|c|c|}
    \hline
    \centering
    \textbf{Prefijo} & \textbf{Símbolo} & \textbf{Factor} & \textbf{Prefijo} & \textbf{Símbolo} & \textbf{Factor}   \\ 
    \thickhline
    deci & d & $10^{-1}$ & deca & da &  $10^{1}$ \\ 
    \hline
    centi & c & $10^{-2}$ & hecto & h &  $10^{2}$  \\ 
    \hline
    mili & m & $10^{-3}$ & kilo & k &  $10^{3}$  \\
    \hline
    micro & $\mu$ & $10^{-6}$ & mega & M &  $10^{6}$  \\
    \hline
    nano & n & $10^{-9}$ & giga & G &  $10^{9}$  \\
    \hline
    pico & p & $10^{-12}$ & tera & T &  $10^{12}$   \\
    \hline
    femto & f & $10^{-15}$ & peta & P &  $10^{15}$  \\
    \hline
    atto & a & $10^{-18}$ & exa & E &  $10^{18}$   \\
    \hline
    \end{tabular}
\end{table}

\normalsize

Para contemplar el espacio entre el valor de la magnitud y la unidad se recomienda usar el comando 
\textcolor{Blue}{\textsf{$\backslash$unit[VALOR]\{UNIDAD\}}}. Por ejemplo, \textcolor{Blue}{\textsf{$\backslash$unit[1,2]\{cm\}}}  da como resultado \unit[1,2]{cm} (no escribir 1,2cm). 

Tenga en cuenta que algunas unidades no están separadas del valor de la magnitud. Ejemplos: 1$\%$, 1º y 1ºC.


\subsection{Formato de las referencias} 
Para las referencias se sigue el formato de las revistas de la APS. 
Citar las referencias utilizando números arábigos consecutivos entre corchetes por orden de aparición en el manuscrito. 

En caso de generar sus archivos en \LaTeX{} se recomienda a los autores que utilicen BibTEX para la generación de las entradas bibliográficas. Si se utiliza BibTEX, durante el proceso de producción final se requerirá el archivo .bib  junto con el archivo .tex principal del manuscrito.

Al usar esta plantilla, el listado de referencias se genera automáticamente por orden de aparición en el texto. 
El comando \color{Blue} \textsf{$\backslash$cite} \color{Black} producira un número de referencia de la forma [\color{blue}NÚMERO\color{Black}]. 

En el listado de referencias, incluir los títulos de los artículos, libros, monografías, tesis, etc. La abreviación de revistas puede consultarse en \urlstyle{same}\url{https://images.webofknowledge.com/WOK46P9/help/WOS/A_abrvjt.html}. 


Datos bibliográficos necesarios para la elaboración de la referencia:

\begin{itemize}[itemsep=1mm, parsep=0pt]
    \item Libros (@book): Autores. Título del libro. Editorial. Lugar de edición. Nro. de edición. Año de publicación. Ver ejemplo en \cite{libro1}.
    \item Artículos de revistas periódicas (@article): Autores. Título del trabajo. Nombre de la revista. Volumen. Número. Páginas. Año de publicación. Ver ejemplos en \cite{paper1} y \cite{paper2}.
    \item Artículos de comunicaciones a congresos (@inproceedings): Autores. Título del trabajo. Título del libro. Páginas donde aparece publicada la contribución. Lugar de edición. Fecha de publicación. Editorial. Ver ejemplo en \cite{proceeding1}.
    \item Tesis doctorales (@phdthesis): Autor. Título de la tesis. Facultad, Universidad. Año. Ver ejemplo en \cite{tesis1}.
    \item Tesis de maestría/licenciatura (@mastersthesis): Autor. Título del trabajo. Facultad, Universidad. Año. Ver ejemplo en \cite{maestria1}.
    \item Material variado (@misc): Título del documento. Dirección electrónica (en caso que corresponda). Fecha de consulta. Opcional: autor. Ver ejemplo en \cite{motor}.
    \item Informe publicado por alguna institución (@techreport): Autores. Título del informe. Institución. Año. Ver ejemplo en \cite{informe1}.
\end{itemize}


%\subsection{Material ilustrativo}
% Las URL deben especificarse así: \urlstyle{same}\url{www.patagonia.com.ar}) 

%-------------------------------------------------------------
% BIBLIOGRAFÍA (edición: AUTORES. Ir a referencias.bib)
\section*{REFERENCIAS} 
 \printbibliography[heading=none] 

%-------------------------------------------------------------
% Fin del documento
\end{document}

